\chapter{测度空间}

\begin{solution}{sigma-algebra-intersections}
设 $(\mathcal A_i)_{i\in I}$ 均为 $X$ 上的 $\sigma$-代数。若 $I=\varnothing$，
相对于 $\mathcal P(X)$ 的空交就是 $\mathcal P(X)$，结论成立。
以下设 $I\ne\varnothing$。$X$ 属于每个
$\mathcal A_i$，故属于其交；若 $E$ 属于其交，则 $E^c$ 属于每个
$\mathcal A_i$；若每个 $E_n$ 属于其交，则 $\bigcup_nE_n$ 属于每个
$\mathcal A_i$。所以交仍为 $\sigma$-代数。

并集一般失败。取 $X=\{1,2,3,4\}$，令 $\mathcal A$ 由分割
$\{\{1,2\},\{3,4\}\}$ 生成，$\mathcal B$ 由
$\{\{1,3\},\{2,4\}\}$ 生成。则 $\{1,2\}\in\mathcal A$、
$\{1,3\}\in\mathcal B$，但它们的交 $\{1\}$ 不属于
$\mathcal A\cup\mathcal B$。
\end{solution}

\begin{solution}{finite-generator-atoms}
若 $m=0$，唯一的交块为 $X$，生成的集合族为 $\{\varnothing,X\}$。
以下设 $m\ge1$。对每个符号串 $\varepsilon\in\{0,1\}^m$，令
$A_\varepsilon=\bigcap_jF_j$，其中 $F_j=E_j$ 当 $\varepsilon_j=1$，
否则 $F_j=E_j^c$。非空的 $A_\varepsilon$ 两两不交并覆盖 $X$。所有这些
原子的并构成一个 $\sigma$-代数，并包含每个 $E_j$，故包含
$\sigma(\mathcal E)$；反过来每个原子由有限交和补集得到，所以属于
$\sigma(\mathcal E)$。因此生成代数正是这些非空交块的任意并。
若其中一个交块含有生成代数的非空子集，该子集也是交块的并；交块两两不交，
故这个子集只能是整个交块。这验证了原子性，原子数至多 $2^m$。
\end{solution}

\begin{solution}{measure-finite-additivity}
对两两不交的 $E_1,\ldots,E_m$，在其后补空集。可数可加性给出
\[
 \mu\Bigl(\bigcup_{k=1}^mE_k\Bigr)=\sum_{k=1}^m\mu(E_k).
\]
若 $E\subseteq F$，应用于
$F=E\mathbin{\dot\cup}(F\setminus E)$ 得单调性。一般有限并依次删去
前面各项化为不交并；再用单调性得到
\[
 \mu\Bigl(\bigcup_{k=1}^mE_k\Bigr)\le\sum_{k=1}^m\mu(E_k).
\]
\end{solution}

\begin{solution}{continuity-from-below-proof}
令 $F_1=E_1$，$F_n=E_n\setminus E_{n-1}$。则 $F_n$ 两两不交，
$E_N=\bigcup_{n\le N}F_n$，而 $E=\bigcup_nF_n$。于是
$\mu(E_N)=\sum_{n\le N}\mu(F_n)$，其极限按非负级数的定义等于
$\sum_n\mu(F_n)=\mu(E)$。
\end{solution}

\begin{solution}{continuity-from-above-counterexample}
在 $(\N,\mathcal P(\N),\#)$ 上取 $E_n=\{n,n+1,\ldots\}$。它们下降到
空集，但 $\#E_n=+\infty$ 对所有 $n$ 成立，不能趋于
$\#\varnothing=0$。失败恰好来自起始集合测度无限。
\end{solution}

\begin{solution}{weighted-counting-measure}
沿用正文的无重复枚举 $X=\{x_k:k\in I\}$ 及 $w_k\in[0,+\infty]$。
空集对应空和，故测度为零。对两两不交的 $(E_n)_{n\ge1}$，每个指标至多
属于一个指标集 $J_n=\{k\in I:x_k\in E_n\}$，且这些指标集两两不交。
记 $J=\bigcup_nJ_n$。则
\[
 \sum_{k\in J}w_k=\sum_{n=1}^{\infty}\sum_{k\in J_n}w_k.
\]
为直接核对这次分组求和，左侧任一有限子和只涉及有限多个 $J_n$，故不超过
右侧。反过来，在有限多个 $J_n$ 中各取有限子集，将它们合并，所得和不超过
左侧；先对这些有限子集取上确界，再对所取组数取上确界，得到反向不等式。
无穷权重也适用这一论证，不需要积分或双重积分定理。因此公式定义测度。

它有限当且仅当 $\sum_{k\in I}w_k<\infty$；它是 $\sigma$-有限的当且仅当
每个 $w_k$ 有限。充分性由单点集覆盖得到，有限 $X$ 时在有限覆盖后补空集；
必要性则因一旦 $w_k=+\infty$，任何包含 $x_k$ 的集合测度均为无穷，
所以不可能用有限测度集合覆盖该点。
\end{solution}

\begin{solution}{dirac-integral-preview}
$\delta_a(\varnothing)=0$。对不交集列，$a$ 至多属于其中一项，所以两边
都等于 $0$ 或 $1$，可数可加性成立。零测集都不含 $a$，其任意子集仍不含
$a$ 且在幂集 $\sigma$-代数中可测，故完备。最后
$\delta_a(E_n)=\mathbf1_{E_n}(a)$，结论只是这一恒等式的改写。

幂集假设不可任意略去。例如在 $X=\{a,b,c\}$ 上，取
$\mathcal A=\{\varnothing,\{a\},\{b,c\},X\}$。限制到 $\mathcal A$
的 $\delta_a$ 仍是测度，但可测零集 $\{b,c\}$ 的子集 $\{b\}$ 不可测，
故这一限制测度不完备。
\end{solution}

\begin{solution}{completion-equivalent-form}
若 $A=E\cup Z$，其中 $Z\subseteq N$、$\mu(N)=0$，则
$E\setminus N\subseteq A\subseteq E\cup N$，两端可测且差集含于 $N$。
反之，若 $E\subseteq A\subseteq F$ 且 $\mu(F\setminus E)=0$，则
$A=E\cup(A\setminus E)$，而 $A\setminus E\subseteq F\setminus E$，
所以 $A$ 属于完备化。
\end{solution}

\begin{solution}{completion-minimality}
设 $(X,\mathcal C,\nu)$ 是包含 $\mathcal A$ 的完备延拓。若
$A=E\cup Z\in\overline{\mathcal A}$，取可测零测集 $N$ 使 $Z\subseteq N$。
因为 $\nu(N)=\mu(N)=0$ 且 $\nu$ 完备，$Z\in\mathcal C$，从而
$A\in\mathcal C$。故 $\overline{\mathcal A}\subseteq\mathcal C$；并且
$\nu(A)=\nu(E)=\mu(E)=\overline\mu(A)$。
\end{solution}

\begin{solution}{measure-algebra-metric}
若 $E\sim E'$、$F\sim F'$，则
$(E\triangle F)\triangle(E'\triangle F')$ 含于
$(E\triangle E')\cup(F\triangle F')$，故测度为零，$d$ 良定义。对称性和
非负性直接成立；由 $\mu(X)<\infty$，距离处处有限。
$d([E],[F])=0$ 正是 $[E]=[F]$。由集合包含
\[
 E\triangle G\subseteq(E\triangle F)\cup(F\triangle G)
\]
及次可加性得到三角不等式。
\end{solution}

\begin{solution}{indicator-set-limits}
固定 $x$。数列 $\mathbf1_{E_n}(x)$ 只取 $0,1$，它收敛当且仅当最终恒定。
最终恒为一等价于 $x\in\liminf E_n$；无穷多次为一等价于
$x\in\limsup E_n$。因此对每个 $x$ 收敛恰好等价于两集合相等，极限就是
共同集合的指标函数。
\end{solution}

\begin{solution}{set-fatou-inequality}
置 $F_N=\bigcap_{n\ge N}E_n$。则
\[
 F_N\uparrow\liminf E_n,
 \qquad
 \mu(F_N)\le\inf_{n\ge N}\mu(E_n).
\]
由上升连续性，
\[
 \mu(\liminf E_n)=\lim_N\mu(F_N)
 \le\lim_N\inf_{n\ge N}\mu(E_n)
 =\liminf_n\mu(E_n).
\]
\end{solution}

\begin{solution}{first-borel-cantelli}
属于无穷多个 $E_n$ 的点集正是 $\limsup E_n$。对每个 $N$，
\[
 \mu(\limsup E_n)
 \le\mu\Bigl(\bigcup_{n\ge N}E_n\Bigr)
 \le\sum_{n\ge N}\mu(E_n).
\]
级数尾和趋于零，故该集合零测。
\end{solution}

\begin{solution}{pushforward-measure}
$T_\#\mu(\varnothing)=0$。若 $B_n$ 两两不交，则其原像也两两不交，且
$T^{-1}(\bigcup_nB_n)=\bigcup_nT^{-1}(B_n)$。于是
\[
 T_\#\mu\Bigl(\bigcup_nB_n\Bigr)
 =\mu\Bigl(\bigcup_nT^{-1}(B_n)\Bigr)
 =\sum_nT_\#\mu(B_n).
\]
故 $T_\#\mu$ 是测度。
\end{solution}
